% test-015.tex -- comandos \spPoly, \spAfig y \spPfig
%
% Compilar:  lualatex-dev test-015.tex
% Validar:   show-pdf-tags --xml test-015.pdf | rnv-wrapp
%            verapdf --flavour ua2 --format text test-015.pdf
%
% Cada caso es un parrafo numerado, para poder referirse a el al escuchar
% el PDF. Los comentarios "doc:" indican lo que documenta el manual; en el
% resto, la lectura se revisa escuchando. Todos trabajan en modo
% matematico. \spPoly recibe tres o mas puntos. \spAfig y \spPfig reciben,
% entre parentesis, un numero (o n), tres o mas puntos, o una de las
% letras F, P, B, L y T con subindice opcional. Restricciones del manual:
% read-arg solo se usa con puntos, y print-sym=triangle o square solo con
% numeros, nunca con F, P, B, L o T. Las entradas invalidas (que detienen
% la compilacion) no van aqui, entre ellas dos puntos o uno solo.
%
% Discrepancia para decidir: el manual de \spPoly dice que con cinco o mas
% puntos no se anade concepto (se lee "A B C D E"), pero el codigo agrega
% el nombre del poligono (pentagono, hexagono, ...). Los casos 3 y 4 sirven
% para escucharlo.
\DocumentMetadata
  {
    lang=es-CL,
    pdfversion=2.0,
    pdfstandard=ua-2,
    tagging=on,
    tagging-setup={math/setup={mathml-SE}},
  }
\documentclass{article}
\usepackage[spanish]{babel}
\usepackage{unicode-math}
\usepackage{spintent}
\usepackage{hyperref}
\hypersetup
  {
    pdfdisplaydoctitle = true,
    pdftitle           = {test-015: spPoly, spAfig y spPfig},
  }
\begin{document}

\section*{A. Comando spPoly: argumentos}

% doc: imprime \triangle ABC -- «triángulo A B C»
1. Tres puntos: $\spPoly{ABC}$.

% doc: imprime \mdlgwhtsquare ABCD -- «cuadrado A B C D»
2. Cuatro puntos: $\spPoly{ABCD}$.

% doc: el manual dice «A B C D E»; el código agrega «pentágono»
3. Cinco puntos: $\spPoly{ABCDE}$.

% doc: el código agrega «hexágono»
4. Seis puntos: $\spPoly{ABCDEF}$.

5. Tres puntos con comas: $\spPoly{A,B,C}$.

6. Tres puntos con comas y espacios: $\spPoly{A, B, C}$.

7. Con primas: $\spPoly{A'B'C'}$.

8. Con subíndices: $\spPoly{A_{1}B_{2}C_{3}}$.

9. Cuatro puntos con subíndices: $\spPoly{A_{1}B_{2}C_{3}D_{4}}$.

10. Con otras letras: $\spPoly{PQR}$.

11. Cuatro puntos con otras letras: $\spPoly{MNOP}$.

\section*{B. Comando spPoly: llaves}

% doc: triángulo A B C, sin el prefijo mayúscula
12. Con read-cmd school, el valor por defecto: $\spPoly[read-cmd=school]{ABC}$.

% doc: MathCAT aplica sus propias reglas
13. Con read-cmd mathcat: $\spPoly[read-cmd=mathcat]{ABC}$.

% doc: el prefijo se verbaliza antes del concepto
14. Con prefix y tres puntos: $\spPoly[prefix={el}]{ABC}$.

% doc: el prefijo se verbaliza antes del concepto
15. Con prefix y cuatro puntos: $\spPoly[prefix={el}]{ABCD}$.

% doc: el prefijo se verbaliza antes del concepto
16. Con prefix y cinco puntos: $\spPoly[prefix={el}]{ABCDE}$.

17. Con concept true, el valor por defecto: $\spPoly[concept=true]{ABC}$.

% doc: solo se verbaliza el argumento
18. Con concept false y tres puntos: $\spPoly[concept=false]{ABC}$.

% doc: solo se verbaliza el argumento
19. Con concept false y cuatro puntos: $\spPoly[concept=false]{ABCD}$.

% doc: solo se verbaliza el argumento
20. Con concept false y cinco puntos: $\spPoly[concept=false]{ABCDE}$.

% doc: imprime \triangle ABC
21. Con print-sym true, el valor por defecto: $\spPoly[print-sym=true]{ABC}$.

% doc: imprime \mdlgwhtsquare ABCD -- «cuadrado A B C D»
22. Con print-sym true y cuatro puntos: $\spPoly[print-sym=true]{ABCD}$.

% doc: imprime ABC; la lectura es la misma
23. Con print-sym false y tres puntos: $\spPoly[print-sym=false]{ABC}$.

24. Con print-sym false y cuatro puntos: $\spPoly[print-sym=false]{ABCD}$.

25. Con silent false, el valor por defecto: $\spPoly[silent=false]{ABC}$.

% doc: no se verbaliza
26. Con silent true y tres puntos: $\spPoly[silent=true]{ABC}$.

% doc: no se verbaliza
27. Con silent true y cuatro puntos: $\spPoly[silent=true]{ABCD}$.

% doc: no se verbaliza
28. Con silent true y cinco puntos: $\spPoly[silent=true]{ABCDE}$.

29. Con todas las llaves: $\spPoly[read-cmd=mathcat,prefix={el},concept=true,print-sym=true,silent=false]{A_{1}B_{2}C_{3}}$.

\section*{C. Comando spAfig: argumentos}

% doc: imprime \mathcal{A} -- «área»
30. Sin argumentos: $\spAfig$.

% doc: imprime \mathcal{A}_1 -- «área uno»
31. Un número: $\spAfig(1)$.

% doc: imprime \mathcal{A}_n -- «área ene»
32. La letra n: $\spAfig(n)$.

33. Otro número: $\spAfig(2)$.

34. Un número de dos cifras: $\spAfig(12)$.

% doc: imprime \mathcal{A}_{\triangle ABC} -- «área del triángulo A B C»
35. Tres puntos: $\spAfig(ABC)$.

% doc: imprime \mathcal{A}_{\mdlgwhtsquare ABCD} -- «área del cuadrado A B C D»
36. Cuatro puntos: $\spAfig(ABCD)$.

37. Cinco puntos: $\spAfig(ABCDE)$.

38. Tres puntos con comas: $\spAfig(A,B,C)$.

39. Con primas: $\spAfig(A'B'C')$.

40. Con subíndices: $\spAfig(A_{1}B_{2}C_{3})$.

% doc: imprime \mathcal{A}_F -- «área de la figura»
41. La figura: $\spAfig(F)$.

% doc: imprime \mathcal{A}_{F_1} -- «área de la figura uno»
42. La figura con subíndice: $\spAfig(F_{1})$.

43. El polígono: $\spAfig(P)$.

44. El polígono con subíndice: $\spAfig(P_{1})$.

45. La parte basal: $\spAfig(B)$.

46. La parte lateral: $\spAfig(L)$.

47. La parte total: $\spAfig(T)$.

48. El total con subíndice literal: $\spAfig(T_{n})$.

\section*{C. Comando spAfig: llaves}

% doc: área del triángulo A B C
49. Con read-cmd school, el valor por defecto: $\spAfig[read-cmd=school](ABC)$.

% doc: MathCAT aplica sus propias reglas
50. Con read-cmd mathcat: $\spAfig[read-cmd=mathcat](ABC)$.

% doc: el prefijo se verbaliza antes de la expresión
51. Con prefix y sin argumentos: $\spAfig[prefix={el}]$.

% doc: el prefijo se verbaliza antes de la expresión
52. Con prefix y un número: $\spAfig[prefix={el}](1)$.

% doc: el prefijo se verbaliza antes de la expresión
53. Con prefix y tres puntos: $\spAfig[prefix={el}](ABC)$.

% doc: el prefijo se verbaliza antes de la expresión
54. Con prefix y la figura: $\spAfig[prefix={el}](F)$.

55. Con concept true, el valor por defecto: $\spAfig[concept=true](ABC)$.

% doc: solo se verbaliza el argumento
56. Con concept false y sin argumentos: $\spAfig[concept=false]$.

% doc: solo se verbaliza el argumento
57. Con concept false y un número: $\spAfig[concept=false](1)$.

% doc: solo se verbaliza el argumento
58. Con concept false y tres puntos: $\spAfig[concept=false](ABC)$.

% doc: solo se verbaliza el argumento
59. Con concept false y cuatro puntos: $\spAfig[concept=false](ABCD)$.

% doc: solo se verbaliza el argumento
60. Con concept false y la figura: $\spAfig[concept=false](F)$.

% doc: solo se verbaliza el argumento
61. Con concept false y el polígono con subíndice: $\spAfig[concept=false](P_{1})$.

% doc: imprime \mathcal{A}_{\square ABCD} -- «área del trapecio A B C D»
62. Con read-arg y cuatro puntos: $\spAfig[read-arg={trapecio}](ABCD)$.

63. Con read-arg y tres puntos: $\spAfig[read-arg={prisma}](ABC)$.

% doc: los espacios pasan a guiones
64. Con read-arg de varias palabras: $\spAfig[read-arg={triángulo isósceles}](ABC)$.

65. Con read-arg y concept false: $\spAfig[concept=false,read-arg={trapecio}](ABCD)$.

66. Con read-arg y prefix: $\spAfig[prefix={el},read-arg={trapecio}](ABCD)$.

% doc: área uno
67. Con read-sub false, el valor por defecto: $\spAfig[read-sub=false](1)$.

% doc: imprime lo mismo -- «área sub uno»
68. Con read-sub true y un número: $\spAfig[read-sub=true](1)$.

% doc: «área de la figura sub uno»
69. Con read-sub true y la figura: $\spAfig[read-sub=true](F_{1})$.

% doc: «área de la figura uno»
70. Con read-sub false y la figura: $\spAfig[read-sub=false](F_{1})$.

71. Con read-sub true y el total: $\spAfig[read-sub=true](T_{n})$.

% doc: imprime \mathcal{A}_{\triangle ABC}
72. Con print-sym auto, el valor por defecto: $\spAfig[print-sym=auto](ABC)$.

% doc: imprime \mathcal{A}_{ABC}; la lectura es la misma
73. Con print-sym none y tres puntos: $\spAfig[print-sym=none](ABC)$.

74. Con print-sym none y cuatro puntos: $\spAfig[print-sym=none](ABCD)$.

% doc: imprime \mathcal{A}_{\triangle_1} -- «área del triángulo uno»
75. Con print-sym triangle y un número: $\spAfig[print-sym=triangle](1)$.

% doc: imprime \mathcal{A}_{\mdlgwhtsquare_1} -- «área del cuadrado uno»
76. Con print-sym square y un número: $\spAfig[print-sym=square](1)$.

77. Con silent false, el valor por defecto: $\spAfig[silent=false](ABC)$.

% doc: no se verbaliza
78. Con silent true y sin argumentos: $\spAfig[silent=true]$.

% doc: no se verbaliza
79. Con silent true y un número: $\spAfig[silent=true](1)$.

% doc: no se verbaliza
80. Con silent true y tres puntos: $\spAfig[silent=true](ABC)$.

% doc: no se verbaliza
81. Con silent true y la figura: $\spAfig[silent=true](F)$.

82. Con todas las llaves: $\spAfig[read-cmd=mathcat,prefix={el},concept=true,read-arg={trapecio},silent=false](ABCD)$.

\section*{D. Comando spPfig: argumentos}

% doc: imprime \mathcal{P} -- «perímetro»
83. Sin argumentos: $\spPfig$.

% doc: imprime \mathcal{P}_1 -- «perímetro uno»
84. Un número: $\spPfig(1)$.

% doc: imprime \mathcal{P}_n -- «perímetro ene»
85. La letra n: $\spPfig(n)$.

86. Otro número: $\spPfig(2)$.

87. Un número de dos cifras: $\spPfig(12)$.

% doc: imprime \mathcal{P}_{\triangle ABC} -- «perímetro del triángulo A B C»
88. Tres puntos: $\spPfig(ABC)$.

% doc: imprime \mathcal{P}_{\mdlgwhtsquare ABCD} -- «perímetro del cuadrado A B C D»
89. Cuatro puntos: $\spPfig(ABCD)$.

90. Cinco puntos: $\spPfig(ABCDE)$.

91. Tres puntos con comas: $\spPfig(A,B,C)$.

92. Con primas: $\spPfig(A'B'C')$.

93. Con subíndices: $\spPfig(A_{1}B_{2}C_{3})$.

% doc: imprime \mathcal{P}_F -- «perímetro de la figura»
94. La figura: $\spPfig(F)$.

% doc: imprime \mathcal{P}_{F_1} -- «perímetro de la figura uno»
95. La figura con subíndice: $\spPfig(F_{1})$.

96. El polígono: $\spPfig(P)$.

97. El polígono con subíndice: $\spPfig(P_{1})$.

98. La parte basal: $\spPfig(B)$.

99. La parte lateral: $\spPfig(L)$.

100. La parte total: $\spPfig(T)$.

101. El total con subíndice literal: $\spPfig(T_{n})$.

\section*{D. Comando spPfig: llaves}

% doc: perímetro del triángulo A B C
102. Con read-cmd school, el valor por defecto: $\spPfig[read-cmd=school](ABC)$.

% doc: MathCAT aplica sus propias reglas
103. Con read-cmd mathcat: $\spPfig[read-cmd=mathcat](ABC)$.

% doc: el prefijo se verbaliza antes de la expresión
104. Con prefix y sin argumentos: $\spPfig[prefix={el}]$.

% doc: el prefijo se verbaliza antes de la expresión
105. Con prefix y un número: $\spPfig[prefix={el}](1)$.

% doc: el prefijo se verbaliza antes de la expresión
106. Con prefix y tres puntos: $\spPfig[prefix={el}](ABC)$.

% doc: el prefijo se verbaliza antes de la expresión
107. Con prefix y la figura: $\spPfig[prefix={el}](F)$.

108. Con concept true, el valor por defecto: $\spPfig[concept=true](ABC)$.

% doc: solo se verbaliza el argumento
109. Con concept false y sin argumentos: $\spPfig[concept=false]$.

% doc: solo se verbaliza el argumento
110. Con concept false y un número: $\spPfig[concept=false](1)$.

% doc: solo se verbaliza el argumento
111. Con concept false y tres puntos: $\spPfig[concept=false](ABC)$.

% doc: solo se verbaliza el argumento
112. Con concept false y cuatro puntos: $\spPfig[concept=false](ABCD)$.

% doc: solo se verbaliza el argumento
113. Con concept false y la figura: $\spPfig[concept=false](F)$.

% doc: solo se verbaliza el argumento
114. Con concept false y el polígono con subíndice: $\spPfig[concept=false](P_{1})$.

% doc: imprime \mathcal{P}_{\square ABCD} -- «perímetro del trapecio A B C D»
115. Con read-arg y cuatro puntos: $\spPfig[read-arg={trapecio}](ABCD)$.

116. Con read-arg y tres puntos: $\spPfig[read-arg={prisma}](ABC)$.

% doc: los espacios pasan a guiones
117. Con read-arg de varias palabras: $\spPfig[read-arg={triángulo isósceles}](ABC)$.

118. Con read-arg y concept false: $\spPfig[concept=false,read-arg={trapecio}](ABCD)$.

119. Con read-arg y prefix: $\spPfig[prefix={el},read-arg={trapecio}](ABCD)$.

% doc: perímetro uno
120. Con read-sub false, el valor por defecto: $\spPfig[read-sub=false](1)$.

% doc: imprime lo mismo -- «perímetro sub uno»
121. Con read-sub true y un número: $\spPfig[read-sub=true](1)$.

% doc: «perímetro de la figura sub uno»
122. Con read-sub true y la figura: $\spPfig[read-sub=true](F_{1})$.

% doc: «perímetro de la figura uno»
123. Con read-sub false y la figura: $\spPfig[read-sub=false](F_{1})$.

124. Con read-sub true y el total: $\spPfig[read-sub=true](T_{n})$.

% doc: imprime \mathcal{P}_{\triangle ABC}
125. Con print-sym auto, el valor por defecto: $\spPfig[print-sym=auto](ABC)$.

% doc: imprime \mathcal{P}_{ABC}; la lectura es la misma
126. Con print-sym none y tres puntos: $\spPfig[print-sym=none](ABC)$.

127. Con print-sym none y cuatro puntos: $\spPfig[print-sym=none](ABCD)$.

% doc: imprime \mathcal{P}_{\triangle_1} -- «perímetro del triángulo uno»
128. Con print-sym triangle y un número: $\spPfig[print-sym=triangle](1)$.

% doc: imprime \mathcal{P}_{\mdlgwhtsquare_1} -- «perímetro del cuadrado uno»
129. Con print-sym square y un número: $\spPfig[print-sym=square](1)$.

130. Con silent false, el valor por defecto: $\spPfig[silent=false](ABC)$.

% doc: no se verbaliza
131. Con silent true y sin argumentos: $\spPfig[silent=true]$.

% doc: no se verbaliza
132. Con silent true y un número: $\spPfig[silent=true](1)$.

% doc: no se verbaliza
133. Con silent true y tres puntos: $\spPfig[silent=true](ABC)$.

% doc: no se verbaliza
134. Con silent true y la figura: $\spPfig[silent=true](F)$.

135. Con todas las llaves: $\spPfig[read-cmd=mathcat,prefix={el},concept=true,read-arg={trapecio},silent=false](ABCD)$.

\section*{E. En fórmulas}

136. Área de un triángulo: $\spAfig(ABC) = \spAfig(1) + \spAfig(2)$.

137. Área total, basal y lateral: $\spAfig(T) = \spAfig(B) + \spAfig(L)$.

138. Perímetro de un cuadrado: $\spPfig(ABCD) = 4 \cdot \spmside{AB}$.

139. Área y perímetro de la misma figura: $\spAfig(ABCD) \quad \spPfig(ABCD)$.

140. Un polígono y su área: $\spPoly{ABC} \Rightarrow \spAfig(ABC)$.

141. En fórmula destacada:
\[ \spPfig(ABC) = \spmside{AB} + \spmside{BC} + \spmside{CA} \].

\end{document}
